\section{Análisis Detallado de Costos del Proyecto (12 meses)}

El éxito financiero y operativo de la plataforma para la denuncia segura depende de una estimación precisa de los recursos. El presente presupuesto se ha personalizado para reflejar los gastos operativos reales del equipo de trabajo, asegurando que la infraestructura soporte los estándares de seguridad y disponibilidad requeridos por la normativa mexicana.

\subsection{Inversión en Hardware (Pago Único)}

La infraestructura de cómputo representa la base técnica del desarrollo. Se han contabilizado los equipos de alto rendimiento necesarios para la ejecución de entornos de microservicios, contenedores (Docker/Podman) y herramientas de análisis estático de código que demandan una carga considerable de memoria RAM y procesamiento.

\begin{table}[h!]
\centering
\caption{Desglose de equipo de cómputo por integrante}
\begin{tabular}{|l|l|r|}
\hline
\textbf{Integrante} & \textbf{Equipo} & \textbf{Costo (MXN)} \\ \hline
Bryan & Acer Nitro (Core i7 8th Gen, 24GB RAM) & \$6,500.00 \\ \hline
Yayo & Huawei Matebook D16 (Core i5-12450H, 16GB RAM) & \$13,500.00 \\ \hline
Bastian & Acer Nitro Lite 16" (Core i5-13420H, 16GB RAM) & \$14,999.00 \\ \hline
\multicolumn{2}{|r|}{\textbf{Total Hardware}} & \textbf{\$34,999.00} \\ \hline
\end{tabular}
\end{table}

\subsection{Gastos Operativos Individuales (Mensuales)}

A diferencia de un cálculo genérico, se aplican los costos específicos reportados por cada miembro del equipo. Estos gastos cubren la conectividad de alta velocidad indispensable para la sincronización de repositorios y la energía eléctrica necesaria para mantener las estaciones de trabajo operativas durante las jornadas de desarrollo y pruebas de integración.

\begin{table}[h!]
\centering
\caption{Desglose de servicios básicos por integrante}
\begin{tabular}{|l|r|r|r|}
\hline
\textbf{Integrante} & \textbf{Luz (Mensual)} & \textbf{Internet (Mensual)} & \textbf{Subtotal Mensual} \\ \hline
Bryan & \$270.00 & \$500.00 & \$770.00 \\ \hline
Yayo & \$250.00 & \$400.00 & \$650.00 \\ \hline
Bastian & \$250.00 & \$389.00 & \$639.00 \\ \hline
\multicolumn{3}{|r|}{\textbf{Total Servicios Mensual (Equipo)}} & \textbf{\$2,059.00} \\ \hline
\multicolumn{3}{|r|}{\textbf{Total Servicios Anual (12 meses)}} & \textbf{\$24,708.00} \\ \hline
\end{tabular}
\end{table}

\subsection{Licenciamiento y Herramientas Especializadas}

Para garantizar la integridad y seguridad de la plataforma, se han seleccionado herramientas líderes en la industria. El uso de \textbf{Termius} permite una gestión cifrada de servidores, mientras que \textbf{Kiuwan} y \textbf{SonarQube} se implementan como filtros críticos de seguridad (SAST) para mitigar vulnerabilidades antes del despliegue. Los costos en moneda extranjera consideran un tipo de cambio base de \$17.34 MXN.

\begin{table}[h!]
\centering
\caption{Cronograma de herramientas y licencias}
\begin{tabular}{|l|c|c|r|}
\hline
\textbf{Herramienta} & \textbf{Costo Unitario} & \textbf{Periodo / Cantidad} & \textbf{Total (MXN)} \\ \hline
Balsamiq & \$237.00/mes & Mes 1 al 3 & \$711.00 \\ \hline
Visual Paradigm & \$123.00/mes & Mes 1 al 6 & \$738.00 \\ \hline
Termius & \$346.80/mes (\$20 USD) & 3 pers. $\times$ 12 meses & \$12,484.80 \\ \hline
VPS & \$389.00/mes & 12 meses & \$4,668.00 \\ \hline
Kiuwan & \$21,675.00/año & Mes 2 al 12 & \$21,675.00 \\ \hline
SonarQube & \$12,484.80/año & Mes 2 al 12 & \$12,484.80 \\ \hline
\multicolumn{3}{|r|}{\textbf{Total Software y Licencias}} & \textbf{\$52,761.60} \\ \hline
\end{tabular}
\end{table}

\subsection{Resumen Consolidado de la Inversión}

El costo total refleja la inversión integral para el desarrollo de la plataforma. La categoría de Capital Humano constituye el rubro principal, contemplando el trabajo especializado de tres ingenieros becarios. Se incluye un margen destinado a la capacitación técnica inicial necesaria para dominar el stack tecnológico seleccionado, asegurando que el desarrollo cumpla con los estándares de calidad exigidos por la institución.

\begin{table}[h!]
\centering
\caption{Resumen consolidado de la inversión (12 meses)}
\begin{tabular}{|l|r|c|}
\hline
\textbf{Categoría} & \textbf{Costo Total (MXN)} & \textbf{Porcentaje (\%)} \\ \hline
Hardware (Inversión Inicial) & \$34,999.00 & 3.48\% \\ \hline
Servicios Individuales (Luz e Internet) & \$24,708.00 & 2.46\% \\ \hline
Licenciamiento y VPS & \$52,761.60 & 5.25\% \\ \hline
Capital Humano y Capacitación & \$892,657.50 & 88.81\% \\ \hline
\textbf{COSTO TOTAL DEL PROYECTO} & \textbf{\$1,005,126.10} & \textbf{100.00\%} \\ \hline
\end{tabular}
\end{table}

\vspace{5mm}
\noindent \textbf{Nota sobre la viabilidad económica:} El 88.81\% de la inversión se concentra en el talento humano, lo cual es característico en proyectos de ingeniería de software de alto impacto social, donde la propiedad intelectual y el diseño de protocolos de seguridad segura son los activos de mayor valor.